\section*{Bab \ref{ch:b1:lab-techniques-1} --- Teknik Laboratorium I: Keselamatan, Pengukuran, dan Pemisahan}

\begin{solution}{exo:b1:lab-techniques-1:1}
\emph{Bahaya}: H225 (cairan sangat mudah menyala, kategori 2) menuntut
\omterm{def:b1:lab-techniques-1:ghs}{kata sinyal} itu, dan label memuat kata yang lebih parah. H225 adalah
\omterm{def:b1:lab-techniques-1:hazard}{bahaya} fisik; H319 (iritasi mata yang serius) dan H336 (kantuk atau
pusing) adalah \omterm{def:b1:lab-techniques-1:hazard}{bahaya} kesehatan. Tindakan pencegahan: tidak ada nyala api
atau percikan api di dekatnya, botol tertutup; kacamata pelindung;
bekerja di tempat yang berventilasi atau di dalam lemari asam.
% ledger: ghs:acetone, hstat:H225, hstat:H319, hstat:H336
\end{solution}

\begin{solution}{exo:b1:lab-techniques-1:2}
(a) \omterm{def:b1:lab-techniques-1:hazard}{Bahaya}, sebuah sifat asam itu. (b) \omterm{def:b1:lab-techniques-1:hazard}{Risiko}, yang kecil karena
paparannya kecil. (c) \omterm{def:b1:lab-techniques-1:hazard}{Bahaya}. (d) \omterm{def:b1:lab-techniques-1:hazard}{Risiko}: \omterm{def:b1:lab-techniques-1:hazard}{bahayanya} sama, paparannya
lebih besar (uap, percikan) dalam keadaan panas dan terbuka.
\end{solution}

\begin{solution}{exo:b1:lab-techniques-1:3}
$R_f = 2.1/6.0 = 0.35$ dan $4.2/6.0 = 0.70$. Noda kedua mempunyai
$R_f$ yang sama dengan pembanding: sampel mungkin mengandungnya (hal itu
konsisten), tetapi $R_f$ yang sama pada satu eluen tidak membuktikan
identitas; eluen kedua, atau penetesan bersama sampel dan pembanding yang
tetap menghasilkan satu noda, memperkuat kesimpulan itu. Noda pada 0.35
pasti merupakan spesi lain.
\end{solution}

\begin{solution}{exo:b1:lab-techniques-1:4}
Setengah lebar \qty{0.05}{mL}, persegi panjang: $u = 0.05/\sqrt3 = \qty{0.029}{mL}$.
Untuk selisih dua pembacaan yang independen,
$u = \sqrt{2} \times 0.029 = \qty{0.041}{mL}$.
\end{solution}

\begin{solution}{exo:b1:lab-techniques-1:5}
Rata-rata \qty{0.25100}{g}; simpangan $+2$, $-2$, $+5$, 0, $-5$ (dalam
\qty{e-4}{g}), jumlah kuadrat \qty{58e-8}{g^2}, $s = \sqrt{58 \times
10^{-8}/4} = \qty{3.8e-4}{g}$; $u = s/\sqrt5 = \qty{1.7e-4}{g}$;
$U = \qty{3.4e-4}{g}$: $m = (0.25100 \pm 0.00034)$~g, $k = 2$.
\end{solution}

\begin{solution}{exo:b1:lab-techniques-1:6}
Satu porsi: $f = 1/(1 + 3 \times 30/50) = 0.357$, \qty{64.3}{\percent}
terekstraksi. Dua porsi \qty{15}{mL}: $f = 1.9^{-2} = 0.277$,
\qty{72.3}{\percent}. Tiga porsi \qty{10}{mL}: $f = 1.6^{-3} = 0.244$,
\qty{75.6}{\percent}.
\end{solution}

\begin{solution}{exo:b1:lab-techniques-1:7}
$z = 0.0012/\sqrt{0.0006^2 + 0.0002^2} = 0.0012/0.00063 = 1.9 \le 2$:
kompatibel, dengan selisih tipis. Dengan \qty{0.0004}{mol/L}: $z = 0.0012/0.00045 = 2.7$,
tidak kompatibel: \omterm{def:b1:titration-methods:titration}{titrasi} yang lebih presisi mengungkapkan perbedaan yang
nyata (sebuah kesalahan dalam pembuatan larutan, atau dalam \omterm{def:b1:titration-methods:titration}{titrasi}).
\end{solution}

\begin{solution}{exo:b1:lab-techniques-1:8}
Q: sedikit larut dalam keadaan dingin, sangat larut dalam keadaan panas.
P hampir tidak melarutkan padatan itu lebih baik dalam keadaan panas
daripada dingin (tidak ada yang mengkristal ketika didinginkan); R
menahan terlalu banyak zat terlarut dalam keadaan dingin (kehilangan
besar). Dengan Q, \qty{5.0}{g} memerlukan sedikitnya
$5.0/12 \times 100 = \qty{42}{mL}$ pelarut mendidih, yang menahan
$0.3 \times 0.417 = \qty{0.13}{g}$ tetap terlarut dalam keadaan dingin:
paling banyak \qty{4.9}{g}, \qty{97.5}{\percent}, yang dapat diperoleh
kembali (sebelum kehilangan apa pun di atas penyaring).
\end{solution}

\begin{solution}{exo:b1:lab-techniques-1:9}
Etanol ($n_D = 1.3611$ pada \qty{20}{\celsius}); air (1.333) dan
sikloheksana (1.4266) tersisih. Indeksnya berubah dengan suhu, sehingga
sebuah nilai tidak berarti apa-apa tanpa suhunya. Nilai 1.350, di antara
air dan etanol, menunjukkan sebuah campuran, misalnya etanol yang
mengandung air.
% ledger: nD:water, nD:ethanol, nD:cyclohexane
\end{solution}

\begin{solution}{exo:b1:lab-techniques-1:10}
Dengan $t = KV/V_a$ dan $n$ porsi, $\ln f_n = -n\ln(1 + t/n)$; ketika
$n \to \infty$, $n \ln(1 + t/n) \to t$ (karena $\ln(1 + \varepsilon) \sim
\varepsilon$), dan $f_n$ menurun menuju $\mathrm{e}^{-t}$ (\cref{prop:b1:lab-techniques-1:multiple-extractions}).
Di sini $t = 3 \times 30/50 = 1.8$, $\mathrm{e}^{-1.8} = 0.165$: paling
banyak \qty{83.5}{\percent}. Tiga porsi sudah memberikan
\qty{75.6}{\percent}, yaitu \qty{91}{\percent} dari limit itu; mencapai
\qty{99}{\percent} darinya memerlukan 32 porsi yang masing-masing kurang
dari \qty{1}{mL}: setelah dua atau tiga porsi, tambahan hasilnya tidak
sebanding dengan kerjanya.
\end{solution}

\begin{solution}{exo:b1:lab-techniques-1:11}
Dengan $t = x - x_0$: $\int_{-a}^{a} (a - |t|)/a^2 \,\mathrm{d}t =
2 \int_0^a (a - t)/a^2 \,\mathrm{d}t = 2 (a^2/2)/a^2 = 1$. Rata-ratanya
adalah
$x_0$ karena simetri; variansnya adalah $2 \int_0^a t^2 (a - t)/a^2
\,\mathrm{d}t = \frac{2}{a^2}\left(\frac{a^4}{3} - \frac{a^4}{4}\right) =
\frac{a^2}{6}$, sehingga $u = a/\sqrt6 = 0.41\,a$, dibandingkan dengan $a/\sqrt3 = 0.58\,a$
untuk kasus persegi panjang: mengetahui bahwa bagian tengah lebih
mungkin mengurangi ketidakpastian.
\end{solution}

\begin{solution}{exo:b1:lab-techniques-1:12}
$c = m/(MV) = 0.25100/(204.2 \times 0.10000) = \qty{0.012292}{mol/L}$.
\omterm{def:b1:lab-techniques-1:uncertainty}{Ketidakpastian relatif}: massa $1.7 \times 10^{-4}/0.2510 = 6.8 \times
10^{-4}$, volume $0.06/100.00 = 6.0 \times 10^{-4}$; gabungan
$\sqrt{(6.8^2 + 6.0^2)} \times 10^{-4} = 9.1 \times 10^{-4}$, yaitu
\qty{0.091}{\percent}; $u(c) = \qty{1.1e-5}{mol/L}$, $U = \qty{2.2e-5}{mol/L}$:
$c = (0.012292 \pm 0.000022)$~mol/L, $k = 2$. Penimbangan sedikit lebih
dominan; kedua sumber itu berukuran sama.
\end{solution}

\begin{solution}{pb:b1:lab-techniques-1:1}
\textbf{1.} GHS02 mudah menyala, GHS05 korosif, GHS06 toksisitas akut,
GHS07 berbahaya atau iritan.
\textbf{2.} Asam salisilat: \emph{Bahaya}, karena H318 (kerusakan mata
yang serius, kategori 1). Aspirin: \emph{Peringatan} (hanya H302).
\textbf{3.} \omterm{def:b1:lab-techniques-1:hazard}{Bahayanya} tetap, paparannya tidak: volume yang kecil, diukur
dan dipakai di dalam lemari asam, dengan sarung tangan, kacamata
pelindung, dan jas laboratorium, botol segera ditutup, tidak ada nyala
api di dekatnya.
\textbf{4.} H314 (luka bakar parah pada kulit dan kerusakan mata):
kacamata pelindung dan sarung tangan. H332 (berbahaya jika terhirup) dan
H226 (cairan dan uap mudah menyala): lemari asam dan tidak adanya nyala
api.
\textbf{5.} Bilas dengan hati-hati memakai air selama beberapa menit,
lepaskan lensa kontak jika memungkinkan, dan lanjutkan membilas
(P305+P351+P338); lalu mintalah saran medis.
\textbf{6.} Ke dalam wadah limbah berair, sesudah dinetralkan, bukan ke
bak cuci.
\textbf{7.} \ce{C7H6O3 + C4H6O3 -> C9H8O4 + C2H4O2}: C $7 + 4 = 9 + 2$,
H $6 + 6 = 8 + 4$, O $3 + 3 = 4 + 2$.
\textbf{8.} $2.00/138.0 = \qty{1.449e-2}{mol}$; paling banyak
$\num{1.449e-2} \times 180.0 = \qty{2.61}{g}$ aspirin.
\textbf{9.} Padatan kasar itu mengandung asam salisilat yang tidak
bereaksi, asam etanoat, dan air: massanya bukan massa aspirin, dan
kemurniannya tidak diketahui.
\textbf{10.} Apa yang tetap terlarut sesudah pendinginan akan hilang:
\omterm{def:b1:precipitation:solubility}{kelarutan} yang kecil dalam keadaan dingin dan besar dalam keadaan panas
berarti produknya larut ketika panas dan kembali hampir seluruhnya
ketika dingin.
\textbf{11.} Setiap mililiter tambahan menahan bagian produknya tetap
terlarut ketika dingin.
\textbf{12.} Pertumbuhan yang lambat menghasilkan \omterm{def:b1:crystals-metals:crystal}{kristal} yang lebih
besar dan lebih teratur, yang memerangkap lebih sedikit cairan induk dan
pengotor; es menurunkan \omterm{def:b1:precipitation:solubility}{kelarutan}, sehingga mengurangi kehilangan.
\textbf{13.} $\qty{0.045}{L} \times \qty{3.3}{g/L} = \qty{0.15}{g}$.
\textbf{14.} $1.95/2.61 = \qty{75}{\percent}$ (\qty{74.7}{\percent}).
\textbf{15.} $807/6 = \qty{134.5}{\celsius}$.
\textbf{16.} Simpangan $-0.5$, $0.5$, $-1.5$, $0.5$, $-0.5$, $1.5$;
jumlah kuadrat 5.5; $s = \sqrt{5.5/5} = \qty{1.05}{\celsius}$.
\textbf{17.} $u_A = 1.05/\sqrt6 = \qty{0.43}{\celsius}$.
\textbf{18.} Setengah lebar \qty{0.5}{\celsius}: $0.5/\sqrt3 = \qty{0.29}{\celsius}$.
\textbf{19.} $1/\sqrt3 = \qty{0.58}{\celsius}$;
$u = \sqrt{0.428^2 + 0.289^2 + 0.577^2} = \sqrt{0.600} = \qty{0.77}{\celsius}$.
\textbf{20.} $U = 2 \times 0.775 = \qty{1.5}{\celsius}$:
$T = (134.5 \pm 1.5)$~\unit{\celsius}, $k = 2$.
\textbf{21.} Produk itu meleleh pada suhu rendah dan dalam rentang
\qty{8}{\celsius}: produk itu tidak murni (asam salisilat, asam etanoat,
air).
\textbf{22.} Karena diberikan sampai satuan, nilai itu terletak dalam $\pm\qty{0.5}{\celsius}$:
$u_{\mathrm{ref}} = 0.5/\sqrt3 = \qty{0.29}{\celsius}$.
\textbf{23.} $2.3/6.0 = 0.38$ dan $3.3/6.0 = 0.55$. Produk kasar
mengandung asam salisilat ($R_f$ sama dengan pembanding) dan aspirin;
sesudah \omterm{def:b1:lab-techniques-1:recrystallisation}{rekristalisasi}, hanya noda aspirin yang tersisa: asam salisilat
telah dihilangkan, setidaknya sampai di bawah batas yang dapat dideteksi
pelat itu.
\textbf{24.} Aspirin terurai ketika meleleh; jika dipanaskan perlahan,
aspirin sempat terurai, dan campuran yang terbentuk meleleh pada suhu
yang lebih rendah. Nilai dalam tabel mengacu pada pemanasan cepat,
seperti pada bangku pemanas.
\textbf{25.} $z = |134.5 - 135|/\sqrt{0.600 + 0.083} = 0.5/0.827 =$
\textbf{0.60}: jauh di bawah 2, titik leleh terukur dan titik leleh
dalam tabel kompatibel.
\end{solution}
